\chapter{Proportionnalité et données}\label{ch:g6:propdata}

Si trois cahiers coûtent $6$ euros, six cahiers coûtent $12$~: le
double de cahiers, le double du prix. Les quantités qui se comportent
ainsi sont \emph{\omterm{def:g6:propdata:def}{proportionnelles}} --- l'une des idées les plus utiles
de toutes les mathématiques, développée plus loin dans
\cref{ch:g7:prop}. Le chapitre se termine par la lecture et le tracé
de graphiques simples.

\section{Quantités proportionnelles}

\begin{definition}[Proportionnalité]\label{def:g6:propdata:def}
Deux quantités sont \emph{proportionnelles}\index{proportionnalité}
lorsque les valeurs de l'une s'obtiennent à partir des valeurs de
l'autre en multipliant toujours par le \emph{même} nombre, appelé
\emph{coefficient de proportionnalité}.
\end{definition}

\begin{example}\label{ex:g6:propdata:table}
Cahiers à $2$ euros chacun~:
\begin{center}
\renewcommand{\arraystretch}{1.3}
\begin{tabular}{l|cccc}
cahiers & $3$ & $5$ & $8$ & $12$ \\
\hline
prix (euros) & $6$ & $10$ & $16$ & $24$ \\
\end{tabular}
\end{center}
Chaque prix est le nombre de cahiers $\times 2$~: le coefficient est
$2$ (le prix unitaire). Un contrôle rapide qu'un tableau est
proportionnel~: tous les «~\omterm{def:g3:division:remainder}{quotients} de colonnes~»
$\frac{6}{3}, \frac{10}{5}, \frac{16}{8}, \frac{24}{12}$ doivent être
égaux --- ici ils valent tous $2$.
\end{example}

\begin{example}[Un contre-exemple]\label{ex:g6:propdata:nonexample}
L'âge et la taille ne sont pas proportionnels~: un enfant de
$12$~ans n'est pas deux fois plus grand qu'un enfant de $6$~ans.
Contrôle numérique~: $1.50$~m à $12$~ans et $1.15$~m à $6$~ans
donnent des \omterm{def:g3:division:remainder}{quotients} $\frac{1.50}{12} = 0.125$ et
$\frac{1.15}{6} \approx 0.19$ --- pas égaux.
\end{example}

\begin{method}[Compléter un tableau de proportionnalité]\label{met:g6:propdata:complete}
Trois façons, à choisir librement~:
\begin{enumerate}
  \item \emph{coefficient}~: trouver le multiplicateur à partir d'une
        colonne complète, l'appliquer aux autres~;
  \item \emph{colonnes}~: une colonne peut être multipliée par un
        nombre, ou deux colonnes additionnées, pour produire une
        nouvelle colonne~;
  \item \emph{retour à l'unité}~: trouver d'abord la valeur pour
        \emph{un} article, puis multiplier.
\end{enumerate}
\end{method}

\begin{example}\label{ex:g6:propdata:methods}
Cinq mugs identiques coûtent $15$ euros~; que coûtent huit mugs~?

\emph{Retour à l'unité}~: un mug coûte $15 \div 5 = 3$ euros, donc
huit mugs coûtent $8 \times 3 = 24$ euros.

\emph{Colonnes}~: $8 = 5 + 3$~; trois mugs coûtent $\frac{3}{5}$ de
$15$, c'est-à-dire $9$~; donc huit mugs coûtent $15 + 9 = 24$ euros.
Même réponse, comme il se doit.
\end{example}

\section{Pourcentages}

\begin{definition}[Pourcentage]\label{def:g6:propdata:percent}
«~$t\,\%$ d'une quantité~» signifie la \omterm{def:g6:fractions:def}{fraction} $\frac{t}{100}$ de
celle-ci~: $25\,\%$ de $60$ est $\frac{25}{100} \times 60 = 15$. Un
pourcentage est une \omterm{def:g6:propdata:def}{proportionnalité} de coefficient $\frac{t}{100}$.
\end{definition}

\begin{example}\label{ex:g6:propdata:percent}
Une classe de $30$ élèves contient $40\,\%$ de filles. Nombre de
filles, étape par étape~: $40\,\%$ signifie $\frac{40}{100} = 0.4$, et
$0.4 \times 30 = 12$ filles. Repères utiles~:
$50\,\%$ est une \omterm{def:g3:division:half}{moitié}, $25\,\%$ un \omterm{def:g3:division:half}{quart}, $10\,\%$ un dixième ---
donc $10\,\%$ de $30$ est $3$, et $40\,\% = 4 \times 10\,\%$ est
$4 \times 3 = 12$~: même réponse, calculée mentalement.
\end{example}

\section{Lire et tracer des graphiques}

\begin{method}[Lire un tableau ou un graphique]\label{met:g6:propdata:read}
Quelle que soit l'image --- tableau, diagramme en barres, graphique en
lignes~:
\begin{enumerate}
  \item lire d'abord le titre et les unités sur les deux axes~;
  \item pour répondre à une question, localiser la bonne barre/colonne,
        puis lire la valeur soigneusement contre la graduation~;
  \item pour comparer, comparer les barres par la hauteur --- mais
        vérifier que l'axe commence à $0$, sinon l'image peut
        tromper.
\end{enumerate}
\end{method}

\begin{example}\label{ex:g6:propdata:chart}
Le diagramme en barres ci-dessous montre le nombre de livres empruntés
à la bibliothèque de l'école pendant une semaine.

\begin{omfigure}
\begin{tikzpicture}
\begin{axis}[
  width=8.6cm, height=5.2cm,
  ybar, bar width=16pt,
  symbolic x coords={lun.,mar.,mer.,jeu.,ven.},
  xtick=data,
  ymin=0, ymax=22,
  ytick={5,10,15,20},
  ylabel={livres empruntés},
  tick label style={font=\small},
  label style={font=\small}]
  \addplot[fill=omDef!30, draw=omDef] coordinates
    {(lun.,12) (mar.,8) (mer.,17) (jeu.,6) (ven.,15)};
\end{axis}
\end{tikzpicture}

{\small Un diagramme en barres~: une barre par jour, hauteur égale au
compte.}
\end{omfigure}

Lectures~: le jour le plus chargé est mercredi ($17$ livres). Mardi et
jeudi ensemble comptent pour $8 + 6 = 14$ livres --- moins que
vendredi seul ($15$). Total de la semaine~:
$12 + 8 + 17 + 6 + 15 = 58$ livres.
\end{example}

\begin{omfigure}
\begin{tikzpicture}
\begin{axis}[omaxis,
  width=8.2cm, height=5.4cm,
  xmin=0, xmax=9.4, ymin=0, ymax=28,
  xlabel={cahiers}, ylabel={prix (euros)},
  xtick={1,...,9}, ytick={4,8,12,16,20,24}]
  \addplot[omDef, thick, domain=0:9] {3*x};
  \addplot[omDef, only marks, mark=*, mark size=1.5pt] coordinates
    {(1,3) (2,6) (3,9) (4,12) (5,15) (6,18) (7,21) (8,24)};
  \draw[dashed, omThm] (axis cs:6,0) -- (axis cs:6,18)
    -- (axis cs:0,18);
\end{axis}
\end{tikzpicture}

{\small Les quantités \omterm{def:g6:propdata:def}{proportionnelles} ont un graphique très
reconnaissable~: les points s'alignent sur une droite passant par
l'origine (ici des cahiers à $3$ euros chacun~; les pointillés lisent
le prix de $6$ cahiers~: $18$ euros).}
\end{omfigure}

\section{Exercices}

\begin{exercise}[$\star$]\label{exo:g6:propdata:1}
Le tableau est-il proportionnel~? Justifier en calculant les
\omterm{def:g3:division:remainder}{quotients}.
\begin{center}
\renewcommand{\arraystretch}{1.2}
\begin{tabular}{l|ccc}
kg de pommes & $2$ & $3$ & $5$ \\
\hline
prix (euros) & $5$ & $7.5$ & $12.5$ \\
\end{tabular}
\qquad
\begin{tabular}{l|ccc}
âge (ans) & $2$ & $4$ & $8$ \\
\hline
taille (cm) & $85$ & $103$ & $130$ \\
\end{tabular}
\end{center}
\end{exercise}

\begin{exercise}[$\star$]\label{exo:g6:propdata:2}
Quatre croissants coûtent $4.80$ euros. Trouver le prix d'un
croissant, puis de sept croissants.
\end{exercise}

\begin{exercise}[$\star$]\label{exo:g6:propdata:3}
Compléter le tableau de \omterm{def:g6:propdata:def}{proportionnalité} (coefficient d'abord~!)~:
\begin{center}
\renewcommand{\arraystretch}{1.2}
\begin{tabular}{l|cccc}
litres de carburant & $10$ & $25$ & $?$ & $60$ \\
\hline
prix (euros) & $18$ & $?$ & $72$ & $?$ \\
\end{tabular}
\end{center}
\end{exercise}

\begin{exercise}[$\star$]\label{exo:g6:propdata:4}
Une voiture consomme $6$~L de carburant pour $100$~km. Combien de
carburant pour $250$~km~? Pour $350$~km~? Quelle distance peut-elle
parcourir avec $27$~L~?
\end{exercise}

\begin{exercise}[$\star$]\label{exo:g6:propdata:5}
Calculer mentalement, en utilisant les repères de
\cref{ex:g6:propdata:percent}~:
$50\,\%$ de $84$~; \quad $25\,\%$ de $200$~; \quad $10\,\%$ de $63$~;
\quad $30\,\%$ de $70$.
\end{exercise}

\begin{exercise}[$\star$]\label{exo:g6:propdata:6}
Dans une école de $450$ élèves, $60\,\%$ mangent à la cantine.
Combien d'élèves est-ce~? Combien n'y mangent pas~?
\end{exercise}

\begin{exercise}[$\star$]\label{exo:g6:propdata:7}
Avec le diagramme en barres de \cref{ex:g6:propdata:chart}~: quels
jours a-t-on emprunté moins de $10$ livres~? Combien de livres de
plus a-t-on empruntés mercredi que jeudi~?
\end{exercise}

\begin{exercise}[$\star$]\label{exo:g6:propdata:8}
Les températures à midi de lundi à vendredi étaient $14$, $16$, $13$,
$17$, $20$ degrés. Dessiner un diagramme en barres de ces données
(choisir une graduation sensée), et lire le jour le plus chaud.
\end{exercise}

\begin{exercise}[$\star\star$]\label{exo:g6:propdata:9}
Une recette pour $6$ personnes demande $450$~g de farine et $3$
œufs. L'adapter pour $10$ personnes. (Pour les œufs, réfléchir avant
d'écrire un \omterm{def:g6:decimals:places}{nombre décimal} d'œufs~!)
\end{exercise}

\begin{exercise}[$\star\star$]\label{exo:g6:propdata:10}
À vitesse constante, une cycliste parcourt $24$~km en $1$ heure.
\begin{enumerate}
  \item Quelle distance parcourt-elle en $2$~h~? En une demi-heure~?
        En $1$~h $30$~min~?
  \item Combien de temps lui faut-il pour $60$~km~?
\end{enumerate}
\end{exercise}

\begin{exercise}[$\star\star$]\label{exo:g6:propdata:11}
Un magasin offre «~$20\,\%$ de réduction sur tout~». Calculer la
réduction et le nouveau prix pour~: un ballon à $15$ euros~; une
raquette à $40$ euros. Le nouveau prix est-il proportionnel à
l'ancien~? Quel est le coefficient~?
\end{exercise}

\begin{exercise}[$\star\star\star$]\label{exo:g6:propdata:12}
Deux bougies de même hauteur sont allumées en même temps. La grosse
se consume entièrement en $6$ heures, la fine en $4$ heures, chacune
à son rythme constant. Au bout de combien de temps la bougie fine
est-elle exactement deux fois moins haute que la grosse~? (Exprimer
les hauteurs restantes après $t$ heures comme \omterm{def:g6:fractions:def}{fractions} de la hauteur
initiale.)
\end{exercise}

\section{Problème~: cartes, plans, et l'art de mentir avec un
diagramme}

\begin{problem}[{Devoir maison --- l'échelle d'une carte est une
proportionnalité~: les longueurs suivent le coefficient, les aires
suivent son carré, et les diagrammes mal tracés trompent l'œil}]
\label{pb:g6:propdata:1}
Toute carte, tout plan, toute maquette obéit à une seule règle~: les
longueurs réelles et les longueurs dessinées sont
\emph{\omterm{def:g6:propdata:def}{proportionnelles}} (\cref{def:g6:propdata:def}). Le coefficient
porte un nom célèbre --- l'\emph{échelle}\index{échelle} --- et il
cache deux pièges que ce problème déclenche exprès~: les \omterm{def:g6:measure:area}{aires} ne
suivent \emph{pas} l'échelle, et les pourcentages, comme les
diagrammes, dépendent entièrement de ce à quoi on les rapporte.

\medskip
\noindent\textbf{Partie I --- Lire une carte.}
Une carte de randonnée est à l'échelle $1:100\,000$~: chaque
centimètre sur la carte représente $100\,000$ centimètres sur le
terrain.
\begin{enumerate}
  \item Convertir $100\,000$ cm en kilomètres. Que représente $1$ cm
        de cette carte, en km~?
  \item Deux villages sont distants de $7.5$ cm sur la carte~; la
        rive d'un lac y forme une courbe de $12$ cm. Donner les
        distances réelles.
  \item Une voie romaine parfaitement rectiligne mesure $20$ km.
        Quelle est sa longueur sur la carte~?
  \item Une seconde carte annonce «~$1$ cm pour $5$ km~». Écrire son
        échelle sous la forme $1: \,?$. Laquelle des deux cartes
        montre le plus de détails pour une même région~?
  \item Dresser un petit tableau (distance sur la carte $1$, $2$,
        $7.5$ cm en regard de la distance réelle) pour la carte de
        randonnée, et expliquer pourquoi une échelle est exactement
        une \omterm{def:g6:propdata:def}{proportionnalité} au sens de la~\cref{def:g6:propdata:def}.
        Quel est le coefficient qui transforme les centimètres sur
        la carte en kilomètres~?
\end{enumerate}

\medskip
\noindent\textbf{Partie II --- Le plan de la chambre de Zoé.}
Zoé dessine le plan de sa chambre --- un rectangle de $4$ m sur
$3$ m --- à l'échelle $1:50$.
\begin{enumerate}[resume]
  \item Quelles sont les dimensions de la pièce sur le plan~?
  \item Son lit mesure $2$ m sur $0.9$ m, et l'ouverture de la porte
        fait $80$ cm de large. Donner les trois mesures sur le plan.
  \item Voici le piège. Calculer l'\omterm{def:g6:measure:area}{aire} réelle de la pièce en
        m$^2$, puis l'\omterm{def:g6:measure:area}{aire} de son plan en cm$^2$. Convertir l'\omterm{def:g6:measure:area}{aire}
        réelle en cm$^2$ (\cref{def:g6:measure:area}) et la diviser
        par l'\omterm{def:g6:measure:area}{aire} du plan~: le \omterm{def:g3:division:remainder}{quotient} vaut-il $50$~?
  \item Expliquer le nombre trouvé~: sur un plan au $1:50$, chaque
        centimètre carré de papier représente un carré réel de
        $50$ cm sur $50$ cm. Combien cela fait-il de cm$^2$ réels~?
        Énoncer la règle~: quand les longueurs sont divisées par
        $50$, les \omterm{def:g6:measure:area}{aires} sont divisées par\,\dots
  \item Un tapis rond couvre $6$ cm$^2$ sur le plan. Quelle \omterm{def:g6:measure:area}{aire}
        réelle couvre-t-il, en cm$^2$ puis en m$^2$~?
\end{enumerate}

\medskip
\noindent\textbf{Partie III --- Comment mentir avec un diagramme (et
avec un pourcentage).}
\begin{enumerate}[resume]
  \item Un magasin a vendu pour $105$ euros le samedi et pour $110$
        le dimanche. Le gérant trace un diagramme en barres dont
        l'axe vertical commence à $100$~: la barre du samedi monte
        de $5$ petites unités, celle du dimanche de $10$. Que
        suggère l'image au sujet du dimanche, et quelle est la
        vérité~? (Calculer la hausse du dimanche en pourcentage du
        samedi, au pour cent près.) Quelle mise en garde de la
        \cref{met:g6:propdata:read} le gérant a-t-il ignorée~?
  \item Retracer (ou décrire) le diagramme honnête, avec un axe
        partant de $0$~: comment les deux barres se comparent-elles
        alors~?
  \item Une collection passe de $40$ à $60$ timbres~: calculer la
        hausse en pourcentage de la valeur de départ. Elle
        redescend ensuite de $60$ à $40$~: calculer la baisse en
        pourcentage de \emph{sa} valeur de départ. Pourquoi les
        deux réponses diffèrent-elles, pour les mêmes $20$ timbres~?
  \item Période de soldes~: un manteau à $100$ euros est affiché à
        «~$-20\,\%$~», et à la caisse on annonce «~$10\,\%$ de plus
        sur le prix soldé~». Calculer le prix final étape par
        étape. La remise totale est-elle de $30\,\%$~? Expliquer au
        client ce qu'elle vaut réellement.
  \item Bouquet final, retour à la carte de la question 4 ($1$ cm
        pour $5$ km)~: une forêt occupe $3$ cm$^2$ de cette carte.
        Quelle est son \omterm{def:g6:measure:area}{aire} réelle, en km$^2$~? Conclure par la
        règle de tout ce problème~: les longueurs suivent
        l'échelle, les \omterm{def:g6:measure:area}{aires} suivent son\,\dots
\end{enumerate}
\end{problem}
